\chapter{Bounded interfaces and independent implementation}
\MARSChapterHeading{Bounded interfaces}
\label{revision:bounded-interfaces}

A parser receives bytes, not trustworthy field values. A packet can stop in the
middle of its header, declare a body that has not arrived, or contain a checksum
that does not agree with its contents. Reading a mapped address successfully does
not establish that the caller supplied that byte. The function must connect every
load to an explicit bound and every output store to an accepted result.

The complete \texttt{bounded-packet.asm} program in \texttt{revision/labs/files}
implements original project \texttt{013-binary-packet-parser}. Its format is
invented for this exercise. It is not an existing network protocol. We will use
its small header to make bounds, byte order, error priority and output preservation
visible in one runnable program.

You should be able to distinguish declared, available and consumed lengths;
justify loads before executing them; decode multibyte fields; and delay output
until all checks succeed. Earlier chapters established these ideas for arrays
and functions. Here the input's own contents influence what the program wants
to read next. The independent exercise replaces the declared body length with
a bounded delimiter search.

\section{State the byte format and the caller's promises}

The function takes an input pointer in \texttt{\$a0}, available count $n$ in
\texttt{\$a1}, and a selected output-record pointer in \texttt{\$a2}. The input
can begin at any byte offset inside a 256-byte arena. The caller guarantees a
nonnegative count no larger than the arena's remaining bytes. The output is a
separate, nonoverlapping, word-aligned 16-byte record, either \texttt{record\_a}
or \texttt{record\_b}. The parser validates the packet within that supplied window;
it does not discover the allocation size from a pointer.

\begin{table}[htbp]
\centering
\small
\begin{tabular}{lll}
\toprule
Byte offset & Field & Rule\\
\midrule
0--1 & Magic & Hexadecimal \texttt{4D 50}\\
2 & Version & One\\
3 & Flags & Only bit zero may be set\\
4--5 & Payload length $L$ & Unsigned, most significant byte first; $0\leq L\leq64$\\
6--7 & Sequence & Unsigned, most significant byte first\\
8 onward & Payload & Exactly $L$ bytes; none when $L=0$\\
$8+L$ & Checksum & Whole packet sum is zero modulo 256\\
\bottomrule
\end{tabular}
\caption{The teaching packet format. All offsets are relative to the passed input pointer. The fixed header is eight bytes; the complete packet consumes $9+L$.}
\label{tab:revision-bounded-format}
\end{table}


\begin{samepage}
Only flag bit zero is accepted; its value is reported as metadata without any
additional behavior. The sequence number is reported as an unsigned integer,
not checked for duplicates or increasing order. Such a policy would need state
across calls and a different interface.\par
\end{samepage}

Length $L$ describes the payload alone. The complete packet requires eight header
bytes, $L$ payload bytes and one trailer, consuming $9+L$. A valid empty payload
still consumes nine bytes. More available bytes may follow the first packet;
those suffix bytes do not contribute to its checksum, output or consumed count.
A zero payload byte is ordinary binary data, not a string terminator.

The default input is hexadecimal
\texttt{4D 50 01 01 00 03 12 34 7F 80 FF 1A}. It declares three payload bytes
and sequence hexadecimal \texttt{1234}, decimal 4660. Its header sum is 232,
payload sum 510 and trailer 26, giving total 768. That total is zero modulo 256.
The accepted output words are flags one, sequence 4660, length three and payload
sum modulo 256, which is 254. Success returns consumed count 12 in
\texttt{\$v0} and payload address $input+8$ in \texttt{\$v1}.

\begin{figure}[htbp]
\centering
\begin{tikzpicture}[x=7mm,y=9mm,alt={Twelve mapped bytes at input offsets zero through eleven are hexadecimal 4D 50 01 01 00 03 12 34 7F 80 FF 1A. The first eight form the header, the next three are payload, and the final byte is the checksum. A bracket for available eight stops after the header and does not permit a trailer read. A second bracket spans all twelve bytes, the complete packet.}]
\fill[black!12] (0,0) rectangle (8,1);
\fill[black!25] (11,0) rectangle (12,1);
\foreach \x in {0,...,12} {\draw (\x,0)--(\x,1);}
\draw (0,0)--(12,0) (0,1)--(12,1);
\foreach \i/\b in {0/4D,1/50,2/01,3/01,4/00,5/03,6/12,7/34,8/7F,9/80,10/FF,11/1A} {
 \node[font=\small\ttfamily] at (\i+.5,.5) {\b};
 \node[font=\footnotesize,anchor=north] at (\i+.5,-.1) {\i};
}
\node[font=\small,anchor=south] at (4,1.2) {Header};
\node[font=\small,anchor=south] at (9.5,1.2) {Payload};
\node[font=\small,anchor=south] at (11.7,1.2) {Trailer};
\draw[|-|] (0,-.7)--(8,-.7);
\node[font=\small,anchor=north] at (4,-.8) {$n=8$: header available};
\draw[|-|] (0,-1.5)--(12,-1.5);
\node[font=\small,anchor=north] at (6,-1.6) {$n=12$: complete packet available};
\end{tikzpicture}
\caption{Mapped bytes and available bytes are different promises. A readable byte outside the passed window remains outside this call's input. The checksum trailer is required even when $L=0$.}
\label{fig:revision-bounded-window}
\end{figure}


These lengths answer different questions. Available tells which bytes the caller
permits this call to read. Declared tells what the header claims the body contains.
Consumed tells how far a successful caller can advance to the next item. Replacing
one with another can accept a truncated packet or consume bytes belonging to its
successor.

\section{Prove availability before interpreting bytes}

The first check is $n\geq8$, before any input load. Even if one available byte
already has the wrong magic, the specified status for fewer than eight bytes is
need-more. The worked entry prepares that error and a zero payload pointer:

\noindent\begin{minipage}{\linewidth}
\MARSCodeLabel{Establish a complete-header bound before the first load}
\begin{minted}{mars-idea}
parse_packet:
    li $v0, -1
    move $v1, $zero
    li $t9, 8
    bltuc $a1, $t9, packet_return
    nop
\end{minted}
\end{minipage}

Once this check passes, every byte at offsets zero through seven is available.
The parser can safely inspect the header, but its values can still be invalid.
Apply the ordered decisions in Table~\ref{tab:revision-bounded-priority}. An
input with multiple defects must receive the first applicable status, so different
correct implementations agree on the result.

\begin{table}[htbp]
  \centering
  \small
  \begin{tabular}{rll}
    \toprule
    Order & Decision                           & Return on failure \\
    \midrule
    1     & Fewer than eight bytes available   & $-1$ (need more)  \\
    2     & Either magic byte differs          & $-2$              \\
    3     & Version is not one                 & $-3$              \\
    4     & Any reserved flag bit is set       & $-4$              \\
    5     & Declared payload length exceeds 64 & $-5$              \\
    6     & Fewer than $9+L$ bytes available   & $-1$ (need more)  \\
    7     & Packet sum is nonzero modulo 256   & $-6$              \\
    \bottomrule
  \end{tabular}
  \caption{Error priority is part of the interface. Apply the decisions from top to bottom; do not let a later defect replace an earlier status.}
  \label{tab:revision-bounded-priority}
\end{table}


For example, eight available bytes with valid magic but a declared length of
65 produce $-5$, not need-more: the semantic length limit is checked before
complete-packet availability. If only seven bytes are supplied, the initial
bound instead gives $-1$ without interpreting that length. These choices are
contract decisions, not consequences of whichever byte the implementation loads
first.

The length stage first decodes the two header bytes, then applies the cap before
forming the required total:

\noindent\begin{minipage}{\linewidth}
\MARSCodeLabel{Separate a permitted length from an available packet}
\begin{minted}{mars-idea}
packet_length:
    lbu $t0, 4($a0)
    sll $t0, $t0, 8
    lbu $t7, 5($a0)
    or $t0, $t0, $t7
    li $v0, -5
    li $t9, 64
    bltuc $t9, $t0, packet_return
    nop
    addiu $t8, $t0, 9
    li $v0, -1
    bltuc $a1, $t8, packet_return
    nop
\end{minted}
\end{minipage}

After the cap, $9+L$ is at most 73, so adding nine cannot overflow a word. Only
after $n\geq9+L$ succeeds may the program read every payload byte and the trailer.
For $L=0$, that second bound still requires nine bytes. The format's small cap
makes this arithmetic proof explicit; a larger length field would need its own
overflow analysis rather than an unchecked total.

Every input access must fit wholly inside the supplied interval. A word load
starting at the last available byte would read beyond it even if its starting
address were legal. Mapped padding can hide such an error from ordinary execution.
The checks therefore inspect native access addresses and widths as well as results.
A valid implementation can reread header bytes; no exact load count is prescribed.

\section{Decode integers without borrowing the machine's byte order}

For high byte $h$ and low byte $l$, a big-endian 16-bit field has value
$256h+l$. The length fragment shifts the first unsigned byte by eight, then ORs
the second into the disjoint low bits. Sequence decoding uses the same rule at
offsets six and seven. Separate byte loads work at every permitted input offset
and make the format's byte order explicit.

For wire bytes \texttt{12 34}, the decoded integer is \texttt{1234}. Storing that
integer as a native word in the default little-endian MARS layout produces bytes
\texttt{34 12 00 00}. This is not a reversed packet field: the output record stores
a decoded integer in a different representation. Do not copy wire bytes directly
into an output word and expect both representations to agree.

Use unsigned loads for these field bytes. With low byte \texttt{80}, a signed
\texttt{lb} yields a negative register value whose high bits are ones. ORing it
with a shifted high byte can set bits outside the intended 16-bit field.
The sequence pair \texttt{FF 80} should decode to 65408, not a negative value.
The permitted sequence range is zero through 65535, independent of its storage
in a signed 32-bit output word.

\begin{samepage}
Pointers have another width obligation. The cursor and returned payload address
must retain the input pointer's full bits. The function may use word arithmetic
for bounded byte counts and sums, but pointer advances use \texttt{daddiu} or
\texttt{daddu}. A short packet does not imply that its address fits in a word.
The input and output pointers are supplied arguments, not instructions to reload
the default global labels.\par
\end{samepage}

\section{Validate the complete sum and retain the payload sum}

The checksum rule is that the sum of header, payload and trailer is zero modulo
256. The worked parser scans the eight header bytes into a total accumulator,
then visits the $L$ payload bytes. Each payload byte contributes both to that total
and to a separate payload accumulator. The trailer contributes only to the total.
Masking the low eight bits tests the total residue and produces the reported
payload residue.

The body loop uses a remaining count. Before a payload iteration, the cursor names
the next payload byte, the count names how many remain, and the two accumulators
cover the bytes already visited. Each iteration loads one byte, updates both sums,
advances the full pointer and decrements the count. When it reaches zero, the
cursor names the trailer. The earlier availability proof establishes that every
one of these accesses lies in the window.

For an empty body, skip the payload loop completely. The initial body cursor
$input+8$ already names the trailer, so executing a first body iteration would
misclassify that checksum byte as payload. After validation, the returned payload
slice still begins at $input+8$, but its length is zero. A pointer alone does not
say how many bytes belong to the slice.

The unsigned total is at most $73\cdot255=18615$ and the payload sum at most
$64\cdot255=16320$, comfortably within a word. Reducing the sums after each byte
or only at the end produces the same residues under these bounds. A payload-only
signed-byte variant can also be correct: each high-bit signed byte differs by
$-256$, which does not change either final residue. That reasoning does not justify
signed loads for sequence decoding, which needs the actual unsigned field value.

An additive checksum detects some changes but cannot distinguish all messages.
Increasing the default payload's first byte from 127 to 128 and decreasing its
second from 128 to 127 preserves both sums. The parser should accept that changed
payload too. The checksum does not authenticate a sender, and its acceptance is
not proof that every field was decoded correctly.

\section{Keep output untouched until acceptance}

All earlier values are tentative. The worked routine keeps flags, sequence,
length and sums in registers until the checksum passes. Only then does it write
the selected output record and prepare successful return values:

\noindent\begin{minipage}{\linewidth}
\MARSCodeLabel{Publish four checked fields and the borrowed payload slice}
\begin{minted}{mars-idea}
packet_commit:
    andi $t4, $t4, 255
    sw $t1, 0($a2)
    sw $t2, 4($a2)
    sw $t0, 8($a2)
    sw $t4, 12($a2)
    move $v0, $t8
    daddiu $v1, $a0, 8
\end{minted}
\end{minipage}

Every failure returns its negative status with \texttt{\$v1}=0 and no output
stores. Leaving the same final bytes is weaker: writing a field and restoring
its old value still violates this contract. Delaying the stores also avoids
needing an output backup or rollback path. The four successful stores are ordinary
sequential operations; this interface makes no concurrent atomic-publication promise.

A debugger at \texttt{packet\_commit} sees a successful validation but can still
see the tentative $-6$ in \texttt{\$v0}. The function has not yet returned.
At \texttt{parse\_ready}, the return interface is complete: the four record words,
consumed count and payload pointer agree. Stop there before main's exit service
overwrites \texttt{\$v0}. Main later stores the status and full payload address
in separate metadata for inspection after normal exit.

The routine is a leaf and uses only caller-clobbered registers. Saved registers
and \texttt{\$sp} remain unchanged. Another implementation may use a properly
restored private stack frame. The caller owns the input bytes backing the returned
slice and must keep them available while using it; the parser made no payload copy.
A later overwrite of that input changes what the borrowed pointer denotes.

On success, a caller can advance its input pointer by the positive consumed count
and subtract that amount from its available count. On need-more it can retain the
same prefix and retry after obtaining more bytes. Other invalid-input statuses
require a caller policy, such as rejecting the record. This parser does not scan
past a malformed prefix to guess where another packet begins.

Tests should exercise successful values, every incomplete prefix, overlapping
error conditions, alternate output records and the absence of forbidden stores.
The original public checks do this without requiring a particular register
allocation or source spelling. Preserve a failing input and locate the first
violated bound, decision or write, rather than adjusting only the final status.

\section{Write a bounded delimiter search independently}

Implement \texttt{read\_line} in \texttt{bounded-line.asm}. It parses one binary
record ending with byte hexadecimal \texttt{0A}, called LF. It accepts at most
32 payload bytes before that delimiter. Other bytes, including zero, \texttt{0D}
and high-bit values, are ordinary payload. There is no escape syntax or text
encoding conversion in this interface.

Arguments are input pointer, nonnegative available count and a selected eight-byte
output record in \texttt{\$a0}, \texttt{\$a1} and \texttt{\$a2}. The window lies
inside a 256-byte arena, at any byte offset, and output is nonoverlapping and word
aligned. Success returns consumed count $L+1$ and the original input pointer in
\texttt{\$v0}/\texttt{\$v1}; output words are $L$ and payload sum modulo 256,
excluding LF. Failure returns $-1$ for need-more or $-2$ for too-long, with
\texttt{\$v1}=0 and no output writes.

Prove each next byte available before loading it. If LF occurs among the first
33 supplied bytes, accept its preceding zero through 32 payload bytes. If 33
bytes are available without LF, the record is too long. With fewer than 33 bytes
and no LF, return need-more. Thus 32 payload bytes without their next byte still
need more input; a following LF succeeds, while a following non-LF proves too-long.
The delimiter test at that boundary must precede rejection for a 33rd payload byte.

Read only through the first LF, or through the first 33 bytes when none has appeared.
Do not inspect trailing records or mapped padding. You may revisit the permitted
prefix, so either a one-pass sum or a delimiter-first, sum-second implementation
is acceptable. Preserve saved registers, the stack pointer, caller storage, input,
guards and the other output record.

The default bytes \texttt{41 00 80 0A 42} consume four bytes, return a length of
three and sum 193, and leave the final \texttt{42} for the caller. Empty input
returns need-more without reading. An immediate LF accepts an empty payload and
consumes one byte. Derive these cases and the exact-capacity boundary before
coding; the written lab supplies public checks and failure investigations.
